\section{Ordinary elimination and absolute isolated zeros}
\label{sec:elimination}

We now turn separation into a finiteness statement for isolated zeros.
Given an escaping sequence of zeros, we construct charts through its
anchors on which the equations vanish and every original free flag
input survives.  A sufficiently small change of a recoverable free
input then gives a different solution near the original one, within the
same retained open domain.  The final step retests the zero formal data on the perturbed sequence;
Theorem~\ref{thm:absolute-zero} states the resulting criterion, including
the ordinary analytic alternative and the required closure operations.

The induction decreases the number of ordinary variables or, at fixed
dimension, the order of a limiting zero.  Its main difficulty is
that restricting to a derivative center need not preserve a zero of
the original field.  Already for $G(y,z)=z^2+c(y)$, the derivative
center $\partial_zG=0$ is $z=0$, but a zero with $z\ne0$ has nonzero
center value $G(y,0)=c(y)$.  We therefore prove the exact-zero
assertion together with two quantitative alternatives for a
center field that need not vanish: one normalizes a value that is at
least a fixed power of a chosen small scale, and the other produces a chart on which the field
is bounded by a prescribed power of that scale.  These alternatives
allow the induction to return from the center to the original anchors.

The return to an anchor must be exact.  The first two subsections give
the estimates for lifting a center chart and correcting its approximate
hit.  Their key requirement is that the powers controlling derivatives,
inverses, and neighborhoods be fixed before the requested approximation
accuracy is increased.  Proposition~\ref{prop:elimination-recursion}
then carries these estimates through the ordinary-variable induction.

Throughout this section the large variables are independent flag
inputs.  In particular, a parameter of an original equation is allowed
to vary when testing whether a zero of the full system is isolated.
Every monomial is evaluated at the \emph{current} physical backgrounds
of the chart.  The regular implicit calculus of
Theorem~\ref{thm:calculus} and the separation theorem are used after
any necessary retest of the primitive packets on an enlarged saturated
flag.

\subsection{Quantitative chart conventions}

Fix a sequence of real anchors, pass to subsequences when making finite
comparisons, and write $(S,t)$ for free large inputs and ordinary inputs
with an interior limiting value.  A subscript $A$ labels the selected
sequence of anchors.  The separately requested record accuracy is
$\mathcal A$; constants and ordinary germ neighborhoods may depend
on $\mathcal A$, whereas power exponents declared uniform do not.
A chart recipe is a finite sequence of
regular ordinary implicit changes, monomial scalings, and the promotions
defined below.  Widths used in these substitutions are current monomials
of strict positive order times bounded regular units with positive
limiting value.  Equations defining its regular implicit changes have
positive series and an invertible exponent-zero Jacobian.  All analytic
evaluations are on bounded arguments with regular limiting values.  The
starting chart and its positive expressions, with their derivatives, are
regarded as given.  Estimates for a recipe concern only the changes it
adds to that starting chart.

An ordinary function will be called \emph{positive} when its two formal
series have no nonzero coefficient at a negative lexicographic exponent;
this terminology does not prescribe its numerical sign.  Its fixed
ordinary derivatives are bounded and converge to their ordinary
exponent-zero germs by Theorem~\ref{thm:calculus}.  Conversely, boundedness
on an anchored real regime and an ordinary neighborhood proves positivity.
We use this converse for normalized quotients; numerical boundedness is
not, by itself, an assertion of analyticity at a zero rate.

For a positive numerical scale $r\to0$, a finite recipe is
\emph{power moderate in $r$} if every inverted monomial $d$ satisfies
\begin{equation}\label{eq:elimination-moderate}
 |d|\ge r^C
\end{equation}
for some fixed $C$, and its finite required forward derivatives, inverse
Jacobian, and inverse neighborhoods at its hits have power bounds in
$r$.  Different finite requests may have different $C$.  Thus a source
ball of radius $c r^N$, an image ball of radius $c' r^{N'}$, and an inverse
Jacobian norm at most $C' r^{-C}$ are permitted.  Constants and starting
points along the sequence may depend on the request.  Comparisons of
monomials are evaluated at the augmented reference anchor or at points
proved sufficiently close to it.  A monomial need not have bounded ratios
at different ordinary inputs in a fixed ordinary box.

There are two kinds of additional coordinates.  Added free scales and
bounded residuals which enter active changes are recorded as outputs
$R$ of the chart.  Temporarily independent bounded residuals will be
called passive records.  An old coordinate is consumed when its
replacement is made; it and its replacement are not both counted as
active recursion variables.  Exogenous inputs $e$ describe a later flag
extension or a pending thickness.  They are held fixed in the recursion,
and their own coordinates are not included among the records $R$ whose
dependence on $e$ will be made small.  This designation is relative
to the pending center recipe; it does not remove an ordinary argument
of the field at the start of a recursive call from the active count.
After a completed lift, every replacement ordinary input on which the
field depends is counted in the next call.  When taking local inverse
Jacobians at fixed original free inputs, new free coordinates are counted
among the coordinates of that inverse, using absolute increments.

\begin{lemma}[Promotion and logarithmic sensitivities]
\label{lem:elimination-promotion}
A nonzero ordinary coordinate $x_A\to0$ can be replaced, after a
saturated flag extension, by
\begin{equation}\label{eq:elimination-promotion}
 x=\varepsilon\rho,\qquad
 \rho=h^\nu e^{-B}>0,\qquad \varepsilon\in\{1,-1\},
\end{equation}
where $h^\nu$ has strictly positive order.  A bounded log remainder
$B$ is retained whenever it is the sole replacement degree of freedom,
even if its values at the anchors are zero.  It may be omitted only
when a free residual already supplies the replacement coordinate.
The replacement has a single new degree
of freedom $\xi$, either a free large variable or a bounded ordinary
residual.  It preserves the old free inputs and does not increase the
number of active ordinary inputs.  With physical ordinary variables and
all other records held fixed, and $M=|\log\rho|$,
\begin{equation}\label{eq:elimination-direct}
 1\le |\partial_\xi\log\rho|\le C(1+M)^C.
\end{equation}
Let $m$ be a fixed exponential monomial satisfying
$\rho^C\le |m|\le\rho^{-C}$ for some fixed $C$ at the reference
anchors.  Every fixed derivative of $\log|m|$ in physical ordinary
backgrounds, bounded records, or absolute free increments has a power
bound in $1+|\log\rho|$ on sufficiently small neighborhoods preserving
the saturated comparisons, using the bounded background jets supplied
by the packet calculus.  These are the independent physical or prefix
coordinates in which those jets are bounded; pullback through a
power-conditioned prefix may introduce its stated power loss in
$\rho$.  In particular,
\begin{equation}\label{eq:elimination-feedback}
 \rho\,\partial_x\log\rho=o(1).
\end{equation}
Equation~\eqref{eq:elimination-promotion}, with current backgrounds
depending on $x$, is therefore a regular implicit equation in $x$.
\end{lemma}

\begin{proof}
Expand $-\log|x_A|$ against the ordered old large scales, subtracting
successive limiting multiples.  The residual is either bounded or one
new independent positive large scale, with coefficient $1$ or $-1$.
Convert only the new residual's chain of logarithms until it meets the
old cuts.  Lemma~\ref{lem:flags-saturation} proves finite termination,
preservation of old free inputs, and the asserted ordinary count.  In
the bounded case take the residual as $B$.  If a new free residual
survives, take its final free member as $\xi$.

At a direct bounded or unconverted residual the derivative in
\eqref{eq:elimination-direct} has absolute value $1$.  Each logarithm
conversion replaces that derivative by its product with the positive
residual scale and a sign.  The other summands are old scales held fixed
for this explicit derivative.  There is only one new chain, so its
derivative has no competing summands.  Every scale encountered is
$O(M)$ at the anchor.  The finite product is therefore at least $1$
and at most a fixed power of $1+M$.

For the sensitivity assertion, a fixed two-sided moderate monomial satisfies
$\log|h^\alpha|=-\sum\alpha_iZ_i$.  Its largest nonzero scale controls
this sum on a saturated ordinary neighborhood.  If it and its inverse
are moderate in $\rho$, that largest scale is $O(1+|\log\rho|)$.
Triangular differentiation of
$\log Z_i=\sum_{j>i}a_{ij}Z_j+b_i$ introduces only finite products of
scales at or below the largest scale involved, and bounded derivatives
of the current backgrounds.  The same assertion follows inductively
for each fixed higher derivative.  A nonnegative-order monomial with
a power-moderate inverse satisfies the required two-sided bounds;
any bounded nonvanishing ordinary factor is handled by its bounded
jets.  This proves the stated bound for the small moderate factors
and their reciprocals.  Absolute
increments of a free scale have derivative $1$ at that scale, and cause
no additional loss in the triangular induction.  Shrinking to an
absolute power-sized box preserves all comparisons.  Equation
\eqref{eq:elimination-feedback} follows because $\rho$ beats every
fixed power of $|\log\rho|$.  The derivative of
$x-\varepsilon\rho(S,t,x)$ in $x$ tends to $1$; its right-hand side is
positive in the series sense.  Theorem~\ref{thm:calculus} applies.
\end{proof}

The logarithmic sensitivity estimate is used only for monomials
satisfying two-sided power bounds at the scale under consideration.  A positive function or
monomial much smaller than every power of that scale is handled as one
positive expression with bounded fixed derivatives, without separating
the large factors in a formal product rule.

\begin{lemma}[A quantitative local inverse estimate]
\label{lem:elimination-local-inverse}
Suppose a map $\Phi$ has, on a ball about $q_0$, bounds
\[
 \|D\Phi(q_0)^{-1}\|\le C r^{-c},\qquad
 \|D^2\Phi\|\le C r^{-b},
\]
and the available source ball has radius at least $c_0r^a$.
There are source and image balls of fixed power radii on which $\Phi$
is invertible.  Their exponents depend only on $a,b,c$, and a
discrepancy $O(r^T)$ within the image ball is corrected by a source
displacement $O(r^{T-c})$.  For inverse derivatives through an order
$k\ge1$, assume in addition power bounds for the forward derivatives
$D^j\Phi$ through order $\max(k,2)$ on that source ball.  Then those
inverse derivatives have fixed power bounds.  A perturbation has the
same conclusions if it is sufficiently power close in
$C^{\max(k,2)}$ on the same available source ball.
\end{lemma}

\begin{proof}
Choose $N>a$ and $N>b+c$.  On a ball of radius $\eta r^N$, for a
sufficiently small fixed $\eta$, the norm of
$D\Phi(q_0)^{-1}(D\Phi(q)-D\Phi(q_0))$ is at most $1/2$.
The map
\[
 q\longmapsto q+D\Phi(q_0)^{-1}(p-\Phi(q))
\]
is a contraction preserving that ball whenever
$|p-\Phi(q_0)|\le c_1r^{N+c}$.  Its fixed point gives the inverse
and the displacement estimate.  Differentiating the inverse equation
expresses each fixed derivative in finite products of the inverse
Jacobian and forward derivatives through the requested order.  The
additional hypotheses give fixed power bounds for these products.
The inverse perturbation identity
$A^{-1}-B^{-1}=A^{-1}(B-A)B^{-1}$ transfers the estimates to a
sufficiently close map.  This argument also supplies explicit slack
for smaller source and image balls.
\end{proof}

We call a local map \emph{power conditioned in a width $r$} if its
requested finite forward derivatives and inverse Jacobian have power
bounds in $r$, it is defined on an available source neighborhood of
power radius, and its restriction there has an inverse on an image
neighborhood of power radius, as in
Lemma~\ref{lem:elimination-local-inverse}.  The coordinates and
absolute free increments are those fixed above.

\subsection{Lifting a finite center recipe}

\begin{lemma}[Taylor replay at current backgrounds]
\label{lem:elimination-lift}
Let an ambient chart have ordinary variables $(y,w)$ near $w=0$.
Suppose a full-dimensional regular recipe on $w=0$ has been constructed
and is power moderate in a width $\rho$, with its records included
among its outputs.  All necessary flag extensions and passive records
may be included before this recipe is used.  Then the recipe can be
lifted to either $w=\rho u$, for bounded ordinary $u$, or a specified
promotion $w=\varepsilon\rho$ with replacement coordinate $\xi$.
At common independent inputs the lifted lower map and record outputs
approximate the center recipe to $O(\rho^T)$ in any specified finite
$C^k$ norm, for arbitrary $T$.

For a thickness substitution, the resulting full map is power
conditioned.  For a promotion the same holds if the inverse center
records at fixed physical $y$ have arbitrarily power-small derivatives
in its exogenous promotion coordinates.  Calibrated approximate hits
whose errors are sufficiently power small can be corrected to exact
hits.  Only finitely many monomial denominator powers are used for
each choice of $T,k$.
\end{lemma}

\begin{proof}
Include every flag node and passive residual needed for the finite
recipe from the start.  An old recipe not using these additions is
revalidated in its original order by the real-bound criterion in
Theorem~\ref{thm:calculus}; old free inputs remain independent.  First
solve the equation for $w$ with its current backgrounds, using
Lemma~\ref{lem:elimination-promotion}.  Replay the old coordinate
replacements subsequently in their chronological order.

Here is the finite algebra used in this replay.  Expand every needed
active-coordinate derivative by the chain rule and every derivative
of a solved coordinate by its implicit Jacobian formula.  The result
uses a finite jet of base positive functions and monomials, previous
solved coordinates, regular Jacobian inverses, and the finite list of
moderate monomial denominators.  For a base positive function $A$
replace its center value by the reverse Taylor polynomial
\begin{equation}\label{eq:elimination-reverse-taylor}
 {\mathcal T}_J A(S,y,w)
   =\sum_{j=0}^{J-1}\frac{(-w)^j}{j!}\,
       \partial_w^j A(S,y,w).
\end{equation}
The derivatives here are in the ambient independent coordinate before
substitution.  On $|w|=O(\rho)$, Taylor's formula with integral remainder
and the bounded derivative calculus gives
${\mathcal T}_JA-A(S,y,0)=O(\rho^J)$; after a prescribed finite
number of independent derivatives and substitutions the loss is only
a fixed power, paid by increasing $J$.

Throughout the replay and the common-domain argument below, first
absorb the reciprocal of every negative-order inverted monomial,
together with all its required jets, as a whole base positive
expression.  This is an exact rearrangement of the expression and
leaves the recipe and its actual map unchanged.  Every divisor
requiring a relative Taylor correction, and every generator retained
for such relative comparisons, may therefore be taken to have
nonnegative order and a power-moderate inverse.  Each satisfies the
two-sided bound in Lemma~\ref{lem:elimination-promotion}.  A zero-order
monomial is constant in the exponential scales; any accompanying
bounded nonvanishing regular factor is handled by its bounded jets.

For such a two-sided moderate divisor $d$, factor its Taylor correction
as the current $d$ times a unit.  To justify the factorization relatively, put
$M=1+|\log\rho|$.  At common other inputs, the logarithmic
sensitivities of $d$ and $\rho$ are $O(M^C)$, and their logarithms
therefore change by $O(\rho M^C)=o(1)$ on the segment from $w$ to
$0$.  A continuity argument, begun at either endpoint, keeps $\rho$
comparable throughout that segment.  Thus
$d(S,y,0)/d(S,y,w)=1+o(1)$.  Taylor remainder estimates divided by
the moderate lower bound give arbitrary power accuracy for this unit
and its required derivatives.  Its inverse is a regular positive
operation.  Bounded monomials without a moderate inverse are treated
as whole base positive functions together with their finite jets.
Passive records occurring in such monomials are independent arguments
until the stage of the replay which consumes them.

This flattening is finite even for a recipe that itself resulted from
earlier lifts: a previous Taylor formula, derivative formula, or unit
inverse is simply another finite operation.  Fix the finite formulas
first, list all denominator losses, and then take the base Taylor
depth larger than their sum and the requested $T,k$.  A corrected
expression approximating a derivative need not be the derivative of
another corrected expression; it has its own formula
\eqref{eq:elimination-reverse-taylor} and differentiated error bound.

Suppose the replay is complete through one stage.  Compare the next
defining field on the two preceding charts at the same independent
inputs, including its unknown.  The flattened formulas have the same
regular ordinary limit and Jacobian, and differ by arbitrarily high
power error.  Their analytic arguments remain in a smaller common
neighborhood of their regular limits.  Real boundedness proves their
positivity before the implicit operation is taken.  At the center
root the new field has an arbitrarily small residual.  Its regular
inverse gives the root with that accuracy, and implicit differentiation
gives the specified differentiated comparison.  Evaluations at these
nearby roots preserve moderate monomial comparisons.  Induction proves
the asserted approximation of the lower map and all its record
outputs.  There is no evaluation of an extracted symbol at a frozen
background in this argument.

For clarity, use the \emph{completed} replay outputs
$y=Y(q,\xi)$, $R=R(q,\xi)$, and $w=W(q,\xi)$, after all
chronological substitutions including the first $w$ equation.
Thus their derivatives already include every dependence through $w$.
At fixed physical $y$, eliminating just the lower block gives the
inverse-record derivative
\begin{equation}\label{eq:elimination-inverse-record}
 K=R_\xi-R_qY_q^{-1}Y_\xi.
\end{equation}
The inverse perturbation identity and the finite derivative bounds
show that the replay changes this expression by $O(\rho^{T-C})$
for a fixed loss $C$ once the center recipe is chosen.  In a thickness
lift the center map does not use $u$, so its inverse records have zero
$u$-derivative.  Hence the replayed inverse records have arbitrarily
power-small $u$-derivatives.  At fixed physical $y$, the equation
$w=\rho u$ consequently has effective derivative comparable to
$\rho$ in $u$: the record contribution is small by logarithmic
sensitivity, and direct current-$w$ feedback is a unit by
\eqref{eq:elimination-feedback}.

For a promotion, regard $L=\log\rho$ as a function of physical
$(y,w)$, records $R$, and the explicit coordinate $\xi$.  The exact
identity $W=\varepsilon\exp L(Y,W,R,\xi)$ holds on the completed
map.  Differentiating it after eliminating $Y_q$ gives
\begin{equation}\label{eq:elimination-width-schur}
 W_\xi-W_qY_q^{-1}Y_\xi
   =\frac{w}{1-wL_w}\,\sigma,
\end{equation}
where the logarithmic Schur factor is exactly
\begin{equation}\label{eq:elimination-schur}
 \sigma=L_\xi+L_R
       (R_\xi-R_qY_q^{-1}Y_\xi)=L_\xi+L_RK.
\end{equation}
The derivatives $L_\xi,L_w,L_R$ in these formulas hold the other
physical variables and records fixed; $K$, by contrast, uses the
total derivatives of the completed outputs.  Thus no term from
record feedback through $w$ has been omitted.  By hypothesis the
center counterpart of $K$ is arbitrarily power
small, and the replay error is $O(\rho^{T-C})$.  The entries of $L_R$
cost only powers of $1+|\log\rho|$, while $|L_\xi|\ge1$.
Thus $|\sigma|\ge1/2$ eventually, and $1-wL_w=1+o(1)$ by
\eqref{eq:elimination-feedback}.  The physical Schur complement
in \eqref{eq:elimination-width-schur} has absolute value at least
$c\rho$.  All other blocks have power bounds.  Block inversion
therefore bounds the \emph{entire} inverse matrix by a power of
$1/\rho$, using only the inverse of $Y_q$ before this step.  In
particular, the smallness of $K$ was proved without assuming the
inverse of the full $(y,w)$ map.

The bounds also hold for any specified finite derivatives on smaller
absolute power-sized boxes.  Such boxes preserve the scale regime;
the finite-order estimates can equivalently be retested on every
sequence of points in them.  Lemma~\ref{lem:elimination-local-inverse}
therefore supplies local image neighborhoods and exact hit correction.

The calibration preceding this correction matters.  For a desired
physical point $(y_A,w_A)$ and reference records $R_A^0$, set
$u=w_A/\rho(S,y_A,w_A,R_A^0)$, or choose the promotion residual by
expanding $-\log|w_A|$ at that same point and those records.  Do not
form this ratio at an imprecisely displaced center.  The lower map and
its records then differ from the reference hit by arbitrary power
errors.  Logarithmic sensitivity makes the resulting error in $w$
arbitrarily power small as well.  The quantitative inverse corrects
both $y$ and $w$ exactly, with unchanged limiting ordinary parameters.
The correction changes records by another fixed power loss, which is
absorbed in the previously chosen accuracy.  Finitely many nested
lifts multiply power losses and remain power moderate.
\end{proof}

\begin{lemma}[Common domains for accuracy refinements]
\label{lem:elimination-common-domain}
Fix a width $r$ equal to a strict positive-order current monomial
times a bounded regular unit with positive limiting value.
All power bounds and moderation in this lemma refer to $r$.
Consider a family of center recipes with a fixed chronological list
of coordinate replacements, a fixed active flag, and fixed scaling
exponents.  Refining its accuracy changes finite Taylor depths and
their finite formulas, but not that list.  Suppose the regular
equations and analytic arguments at each stage have the same
ordinary limiting germs throughout the family.  Suppose also that,
for every chronological prefix, its actual map and record outputs,
the finite forward derivatives required below, its inverse where
used, and its available source neighborhood have power bounds with
exponents fixed before the accuracy request.  Constants and ordinary
germ neighborhoods may depend on accuracy.  The inverted monomials
are, after the preceding whole-reciprocal preprocessing, products of
powers of a fixed finite list of two-sided moderate generators;
the powers in a flattened formula may depend on accuracy.  Here a
fixed monomial generator means a fixed exponent vector evaluated at
the current backgrounds, not a frozen numerical function.

Fix $J_0>0$ and a finite derivative request, and use the new
transverse width $\rho=r^{J_0}$.  The Taylor replay in
Lemma~\ref{lem:elimination-lift} can then be made,
through every chronological stage, on source balls
$B(q_A,c_{\mathcal A} r^N)$ with $N$ independent of the accuracy request $\mathcal A$.
The same holds for graph tubes used in its regular implicit stages,
with fixed power exponents.  On these domains the comparison can be
made arbitrarily power accurate in the requested differentiated norm.
The resulting family has the same prefix properties, including
power-sized image neighborhoods whenever the Schur hypotheses of
Lemma~\ref{lem:elimination-lift} hold.
\end{lemma}

\begin{proof}
We first justify replay on actual ordinary neighborhoods; estimates
on a power tube alone would not justify a positivity test.  A
two-sided moderate generator $d$ satisfies $d\ge r^C$ at the reference
anchors.
Enlarging $C$ to $C'>C$ makes $r^{C'}/d$ a strict positive-order
monomial times a bounded regular unit.  The flag hierarchy is
preserved under independent bounded ordinary variations, so
$d\ge r^{C'}$ eventually on a smaller ordinary neighborhood.
This argument concerns orders at each input.  It does not compare
the numerical values of $d$ or $r$ at different inputs.  The same
assertion for $d^p$ has exponent $pC'$; $p$ may depend on the chosen
finite recipe.

Fix one member of the center family and flatten its old finite
formulas before selecting the new replay Taylor depth.  All their
denominator powers and chain-rule losses are now a finite fixed
list.  The new Taylor depth introduces higher base jets and more
summands, but no new monomial divisor: a base jet is kept as a
whole positive function, and differentiating the new finite sum
to the fixed requested order loses only that order and the
previously listed prefix powers.  Factoring a divisor correction
uses the same old divisor and a regular unit.  Thus the loss to
be beaten does not grow with the new Taylor depth.
The finitely many base jets finally
needed at that request are defined and bounded on some ordinary
box $U_{\mathcal A}$ with radius independent of $r$, by
Theorem~\ref{thm:calculus}.  Its radius and bounds may depend on $\mathcal A$.
Use a smaller box with spare room.  At every input in this box,
the segment of length $O(r^{J_0})$ in the new transverse coordinate
stays in the ambient base box eventually.  The current and center
widths are comparable along that segment by their logarithmic
sensitivities.  The reverse Taylor remainder and the moderate
generator bounds therefore hold on the entire smaller ordinary
box, pointwise in its current width.

Proceed through the chronological stages.  By the preceding stages,
all their solved ordinary outputs are defined on ordinary boxes.
Flatten the next defining field and all its analytic arguments on
that previous chart.  A finite denominator power loses a finite
power of $r$; increase the new Taylor depth after that loss is
known.  Thus each corrected field or argument differs from its
center counterpart by an arbitrarily small power on a genuine
ordinary box.  Its real limit is the center's limiting ordinary
germ.  Analytic argument tuples stay in the domain of evaluation
with spare ordinary room.  The real-bound criterion now proves
positivity on this box, before the next analytic evaluation or
implicit inversion is performed.  The defining equation has the
same regular limiting Jacobian as its center equation.
Theorem~\ref{thm:calculus} supplies the regular branch on a smaller
ordinary box and bounded finite jets of its ordinary outputs.
This proves the stage induction on ordinary boxes.  Their radii
can shrink with the finite request, but no radius in this argument
is a power $r^{\mathcal A}$ depending on the accuracy.  In particular, an
analytic argument that is bounded only on such a shrinking tube
would not pass this construction.

We next choose common power boxes.  There are finitely many
structural stages.  Record the preassigned source exponents and
forward derivative exponents of their center prefixes, taking all
independent coordinates and the unknowns of later regular
equations as parameters when needed.  If an intermediate target
source ball has exponent $a_j$ and the map into it has derivative
loss $b_j$, require $N>a_j+b_j+1$; include $a_j=0$ for an
ordinary box.  There are only finitely many such requirements.
On a source displacement of size $c_{\mathcal A} r^N$,
each center ordinary argument then moves by $O_{\mathcal A}(r^{N-b})$ for
one of these fixed losses $b$, and so remains inside any of its
accuracy-dependent ordinary boxes for sufficiently far anchors.

The same choice, enlarged by a fixed amount if necessary, keeps
every two-sided moderate generator comparable to its reference value.
Indeed in the independent coordinates of a prefix its logarithmic
derivative is bounded by a fixed power of $1+|\log r|$.
Pullback by a prefix costs at most $r^{-b}$ with $b$ already
recorded.  Integration on the source ball gives
\[
 |\Delta\log d|+|\Delta\log r|
 \le C_{\mathcal A} r^{N-b}(1+|\log r|)^C=o(1).
\]
This estimate is continued from the center as long as the ratios
stay, say, between $1/2$ and $2$; it prevents the first exit, and
hence proves comparability on the whole ball.  Replacing $d$ by a
power $d^p$ multiplies this logarithmic estimate by the fixed
number $p$ for that request.  It changes $C_{\mathcal A}$ and the starting
point, not the chosen exponent $N$.  Nonmoderate positive
expressions are never inverted and retain their bounded-jet
estimates as whole functions.  Saturation preserves the remaining
scale comparisons under these small changes.

For completeness, carry out the regular implicit stages on these
preassigned boxes, rather than extracting a new unspecified box
after each solve.  For the next stage let $v_0(q)$ be the center
root.  Choose a trial tube
\[
 |v-v_0(q)|\le c_{\mathcal A} r^E
\]
whose exponent $E$ exceeds, with spare room, the sum of each
preceding target-domain exponent and the derivative loss of the
map into that domain.  Include earlier trial-tube exponents among
these target-domain exponents.  For the contraction also require
$E>b+c+1$ if the normalized equation has second-derivative
loss $b$ and inverse-Jacobian loss $c$.  Choose these trial
exponents in chronological order; every requirement at a new
stage uses exponents already fixed at earlier stages.  These
losses are supplied by the prefix bounds; in the equation's own
ordinary coordinates its limiting Jacobian is a unit and its
fixed jets are bounded.  Ordinary
argument slack can again be absorbed into $c_{\mathcal A}$ and a later start.
On this tube, use the ordinary-box comparison just established,
after pullback by the preceding prefixes.  If the finite formula
at this request loses a power $C_{\mathcal A}'$, choose its Taylor depth so
that the residual and its required derivatives are
$O(r^{T-C_{\mathcal A}'})$.  For a desired output comparison $O(r^P)$,
choose explicitly
\[
 T-C_{\mathcal A}'>\max\{E+c,\ P+H\}+1,
\]
where $H$ bounds the sum of the fixed inverse, derivative, and
composition losses for the requested differentiated output.
Thus even after multiplication by the inverse norm $r^{-c}$
the correction is smaller than the trial radius $r^E$.
At each $q$ the fixed-Jacobian contraction
used in Lemma~\ref{lem:elimination-local-inverse} then gives a
root strictly inside the trial tube.  The field's Jacobian and
its first variation stay within the chosen contraction bounds
there.  Uniqueness makes the roots at neighboring $q$ agree, so
they define a graph on the whole preassigned source ball.
Implicit differentiation gives the requested comparison on that
same ball.

This argument is sequential: before the next stage, the previous
ones are already defined with spare room on its trial inputs.
Increase $N$ and the finitely many trial exponents at the outset
to meet all prefix requirements.  Each requirement involves only
the fixed structural list, width exponents, center derivative
bounds, and finite derivative order.  Higher Taylor depths affect
only constants, the error order to be requested next, and the
far-enough start.  They add no coordinate replacement or new
domain condition.  Thus these power exponents are independent
of record accuracy.

Finally, the differentiated comparison transfers the center
prefix forward bounds to the actual replayed prefixes on these
common boxes.  The regular implicit stages have unit inverses
in their solved directions; width stages use the Schur estimates
of Lemma~\ref{lem:elimination-lift}.  All other inverse blocks have
the preassigned prefix bounds.  The local inverse lemma, now with
its available source radius explicitly supplied and its required
higher forward bounds, produces image radii with fixed exponents.
This proves the family and prefix conclusions.
\end{proof}

\subsection{The simultaneous elimination invariant}

A numerical floor is a sequence $\delta_A>0$ tending to zero; it is
not initially a flag variable or a function on the chart.  After a
subsequence, a numerical value is \emph{visible} if its absolute value
is at least $\delta_A^K$ for some fixed $K>0$, and \emph{negligible}
if it is smaller than every fixed positive power of $\delta_A$.
A floor realization is a saturated extension, preserving active free
inputs, and a positive width $r$ equal to a strict positive monomial
times an optional bounded positive log-remainder unit, such that
\begin{equation}\label{eq:elimination-floor}
 \log(1/r_A)\asymp\log(1/\delta_A)
\end{equation}
at the augmented reference input.  Its additional bounded coordinates
are exogenous, not active recursion variables.

The exact-zero assertion is proved together with two quantitative
alternatives because a restriction to a derivative center need not
remain an exact-zero problem.  The center quantity $H$ introduced in
Equation~\eqref{eq:elimination-H} may be visible or negligible even
when the original field vanishes.  A visible center supplies a
normalization; a negligible center supplies a chart accurate enough
to lift and correct to the required exact hit.  The conditioning
powers below are fixed before record accuracy, so increasing that
accuracy can make the correction smaller than the already available
inverse neighborhood.

\begin{proposition}[Elimination with a deferred floor]
\label{prop:elimination-recursion}
Let $F$ be a real positive scalar field on an active chart with $n$
ordinary inputs.  The following statements hold simultaneously.
\begin{enumerate}
\item In exact mode, along anchors where $F=0$, finitely many regular
changes, promotions, and exact ordinary restrictions produce a chart
through those anchors on which $F$ has identically zero formal data
and hence vanishes on its real regime.  All old free inputs survive.
\item If $F_A$ is visible relative to a numerical floor, there is a
full-dimensional regular chart with an exact hit and
\begin{equation}\label{eq:elimination-visible}
 F=\tau U,
\end{equation}
where $\tau$ is a nonnegative-order monomial, possibly $1$, and $U$
is positive with nonzero exponent-zero value at the limiting hit.
All added divisors, widths, and local conditioning costs are power
moderate in $|F_A|$.  The active ordinary count does not increase.
\item If $F_A$ is negligible, active flag additions and their reference
records can be chosen without realizing the floor.  For every eventual
realization $r$, field tolerance $M$, finite derivative request $k$, and
record accuracy $\mathcal A$, there is a full-dimensional regular chart hitting
the input, of active ordinary count at most $n$, on which
\begin{equation}\label{eq:elimination-coarse}
 |F|\le r^M.
\end{equation}
Its added divisions and widths are power moderate in $r$.  Its records
at the hit differ from the references by $O(r^{\mathcal A})$, and their inverse
first derivatives in the exogenous inputs, at fixed physical inputs,
are $O(r^{\mathcal A})$ on a power-sized neighborhood.
\end{enumerate}
The third assertion has the following stronger order of quantifiers.
For fixed $M,r,k$, the thickness exponents, forward bounds for the
actual map \emph{and its record outputs} through order at least
$\max(k,2)$, inverse-Jacobian powers, and source and physical-image
neighborhood exponents can all be fixed \emph{before} $\mathcal A$ is requested.
Constants and sufficiently far starting points may depend on $\mathcal A$.
The denominator powers in an individual finite flattened formula may
depend on $\mathcal A$.  Completed active charts remain fixed as $\mathcal A$ increases;
only pending thickness recipes and their finite Taylor replays are
refined.

The construction also has the following prefix invariant.  For the
fixed request before choosing $\mathcal A$, there is a fixed finite
chronological list of coordinate replacements and scaling exponents.
Every prefix has the same ordinary limiting regular equations and
analytic argument germs throughout the accuracy family.  Its maps
and records, with the unknown of a subsequent equation still an
independent ordinary input, have preassigned finite derivative and
source-domain powers; its inverse and image powers are preassigned
where an inverse is used.  Its monomial divisors belong to powers
of a fixed finite list of moderate current monomials and their
triangular derivative factors.  Only their multiplicities inside
the finite formulas can grow with accuracy.  Thus the family meets
the hypotheses and conclusions of
Lemma~\ref{lem:elimination-common-domain} at every pending lift.
\end{proposition}

\begin{proof}
We use induction on $n$, and, at fixed $n$, on the nonzero directional
order of the ordinary leading germ at the limiting hit.  Every step is
performed at the exact hit of completed preceding steps.  We verify
the quantitative third assertion, including its prefix invariant,
in the same induction, not after constructing an uncontrolled family
of charts.  For an identity exit the structural list is empty.
For the single thickness exit its one exponent is fixed by the
field tolerance; all ordinary limiting germs are independent of
accuracy and its direct inverse costs that fixed power.

\paragraph{Normalization and immediate exits.}
If both formal data arrays are zero, separation gives $F=0$ on the
real regime.  This proves the exact assertion and the coarse assertion
with the identity chart and no new records.  Otherwise separation
and leading normalization give
\begin{equation}\label{eq:elimination-leading}
 F=aG,\qquad G\longrightarrow f,
\end{equation}
where $a$ is a nonnegative-order monomial and $f$ is a nonzero real
analytic germ.  The normalized expression $G$ has raw packet trees,
and its locally uniform real convergence makes it positive by the
real-bound criterion in Theorem~\ref{thm:calculus}.  That theorem then
gives the differentiated convergence in
\eqref{eq:elimination-leading}.  If $a$ is negligible in floor mode,
the identity chart already gives \eqref{eq:elimination-coarse} for
every fixed $M$ after any realization: the valuation of $a$ exceeds
every floor power and $G$ is bounded.  No division by $a$ is used in
this output.  If $a$ is visible, it is a moderate divisor and
visibility or negligibility is unchanged on passing from $F$ to $G$.
If $f$ is nonzero at the limiting hit, $G$ is a regular unit, proving
the visible assertion; exact zeros are impossible.  These observations
also prove the induction when $n=0$.

\paragraph{A regular derivative coordinate.}
At a zero of $f$, choose an ordinary direction $z$, complementary
coordinates $y$, and $m\ge1$ such that
\[
 \partial_z^jf=0\quad(0\le j<m),\qquad
 \partial_z^mf\ne0
\]
at the limiting point.  This is possible by taking a direction where
the first nonzero homogeneous Taylor term does not vanish.  Define
\begin{equation}\label{eq:elimination-w}
 w=\partial_z^{m-1}G.
\end{equation}
Its $z$ derivative has nonzero limit, so it is a regular ordinary
coordinate with $w_A\to0$.  In exact mode, if $w_A=0$ along a
subsequence, restrict exactly to $w=0$ and apply the smaller-$n$
assertion.  If $m=1$, then $G=w$: a visible nonzero $w$ is promoted,
whereas negligible $w$ is replaced by $r^Ju$ with $J$ selected from
the field tolerance.  The latter hits with $u\to0$, has no new active
records, and has inverse cost $r^{-J}$ independent of record accuracy.
Together with the exact restriction this proves the $m=1$ step.

Assume $m>1$.  On the center $w=0$ let
\begin{equation}\label{eq:elimination-H}
 c_\ell=(\partial_z^\ell G)|_{w=0}\quad(0\le\ell\le m-2),
 \qquad
 H=\sum_{\ell=0}^{m-2}c_\ell^{2L/(m-\ell)},
\end{equation}
where $L$ is a positive common multiple of $1,\ldots,m$.
The exponents in this sum are even integers; hence $H\ge0$ and
its smallness controls each $c_\ell$.  At the center projection with
the same $(S,y_A)$ these coefficients tend to zero.  Apply the
smaller-$n$ assertion to $H$.  Its numerical floor is
\begin{equation}\label{eq:elimination-theta}
 \theta_A=|w_A|
 \quad\hbox{if $w_A\ne0$ is visible, or in exact mode};
 \qquad
 \theta_A=\delta_A\quad\hbox{otherwise}.
\end{equation}
In the latter case $w_A$ is negligible, with $w_A=0$ allowed.

\paragraph{A visible center and bounded $w/s$.}
Suppose $H$ is visible in its center call.  Write its visible output
as $H=s^{2L}U_H$, with $U_H$ a positive numerical unit and $s$ a
strict positive monomial.  Taking a fixed real power of the output
monomial is allowed.  Then $s_A\ge\theta_A^K$ for some $K$, all
center costs are moderate in $s$, and
\begin{equation}\label{eq:elimination-c-bounds}
 |c_\ell|\le C s^{m-\ell}.
\end{equation}
The quotients in this inequality are positive by real boundedness.

If $w_A/s_A$ stays bounded, use Lemma~\ref{lem:elimination-lift}
with $w=su$.  Choose replay accuracy larger than every power needed
below and correct its hit exactly.  At common lower parameters, the
lifted $y$ and records differ from their center values by arbitrary
power errors, and the current and center $s$ have ratio tending to
$1$.  Taylor expansion in the original independent $z$ at fixed
current $(S,y)$ gives
\[
 z-z_0=b su+O(s^2),\qquad
 b=(\partial_z^mG(S,y,z_0))^{-1},\qquad w(S,y,z_0)=0.
\]
The coefficient of $(z-z_0)^{m-1}$ in $G$ is exactly zero.  Dividing
the Taylor expansion by $s^m$ and using
\eqref{eq:elimination-c-bounds} shows uniform boundedness, and gives
the ordinary limit
\begin{equation}\label{eq:elimination-polynomial}
 \frac{G}{s^m}\longrightarrow
 P(u)=\sum_{\ell=0}^{m-2}\frac{\gamma_\ell b_0^\ell}{\ell!}u^\ell
       +\frac{g_m b_0^m}{m!}u^m,
 \qquad
 \gamma_\ell=\lim\frac{c_\ell}{s^{m-\ell}}.
\end{equation}
Here $g_m=\lim\partial_z^mG\ne0$ and $b_0=g_m^{-1}$.
The coefficients are ordinary germs in the remaining limiting
parameters and are independent of $u$.  Errors caused by replacing
$z-z_0$ with $bsu$ are $O(s)$ after normalization, by
\eqref{eq:elimination-c-bounds}.  At least one $\gamma_\ell$ is
nonzero at the limiting remaining parameter, since $H/s^{2L}$ has
nonzero limit.  Real boundedness proves
that $G/s^m$ is positive with exponent-zero germ $P$.

Every zero of $P$ has multiplicity less than $m$.  Indeed an $m$-fold
root $u_0$ would give $P(u)=c(u-u_0)^m$.  Its absent degree-$m-1$
coefficient forces $u_0=0$, which would make all lower coefficients
zero, a contradiction.  At a limiting nonzero value take the unit
exit.  Otherwise apply the same-$n$ induction at the smaller
directional order in $u$.  In floor mode the factor $as^m$ is visible,
so the inherited numerical floor suffices.  If the final output is
visible, all preliminary costs are moderate in $|F_A|$, because
$|F_A|\le C a_As_A^m$ and both $a,s$ are bounded above.  If the
final output is negligible, these completed active changes have
floor-moderate costs and are fixed before the later accuracy request.
All their chronological prefixes are fixed recipes.  Their ordinary
limiting germs and analytic argument slack are consequently fixed,
and they have a finite list of monomial generators and finite
derivative and domain powers in any later floor realization.
Compose the smaller-order family's prefix bounds with these fixed
prefixes.  A source displacement $r^N$ stays in each required
intermediate source box by taking $N$ larger than the sum of its
fixed radius exponent and the preceding derivative losses.
This choice is independent of the subsequent accuracy.  The
inverse derivative formulas and the local inverse lemma give
the corresponding fixed inverse and image powers.  This proves
preservation of the prefix invariant through the smaller-order
branch; it does not require rebuilding an active change when
the pending accuracy is increased.

\paragraph{A visible center and unbounded $|w|/s$.}
If $|w_A|/s_A\to\infty$, the center floor in
\eqref{eq:elimination-theta} must be $|w_A|$: a negligible $w_A$
could not dominate visible $s_A$.  Promote $w$ against the center's
already extended flag, as an exogenous addition to that recipe,
calibrated at the desired physical point and reference records.  Put
$\rho=|w|$.  The center recipe is independent of the new coordinate,
$s\ge\rho^K$, and $s/\rho\to0$.  Lift it by
Lemma~\ref{lem:elimination-lift}.  The same Taylor formula now gives
\begin{equation}\label{eq:elimination-large-w}
 \frac{G}{\rho^m}\longrightarrow
 \frac{g_m b_0^m\varepsilon^m}{m!}\ne0.
\end{equation}
Indeed each lower term is bounded by
$C(s/\rho)^{m-\ell}$ after normalization.  This is the visible
output with moderate conditioning.  It excludes exact zeros in this
case.

\paragraph{A negligible center relative to nonzero $w$.}
Suppose $H$ is negligible and its floor is $\theta_A=|w_A|$.
First finish the center's active flag additions and select their
reference records.  Then realize this one pending floor by promoting
$w$, expanding $-\log|w_A|$ at the desired physical input and the
reference records, and converting only its new residual chain.
The result is a permitted realization $\rho=|w|$: changing a bounded
background from that physical point to its center projection costs
$O(|w_A|)$, so the logarithmic comparisons of all moderate scales are
unchanged, by Lemma~\ref{lem:elimination-promotion}.

Request a center field tolerance sufficiently large that every
$c_\ell$ will be $o(\rho^{m-\ell})$, including the losses of the
lift.  Fix its map, record-output, inverse, and neighborhood powers
using the strengthened smaller-$n$ assertion at that tolerance and
realization.  Only now request reference-record and inverse-record
accuracy beyond those powers.  Choose the Taylor replay even more
accurately.  After eliminating the lower physical $y$ block, the
logarithmic promotion Schur factor is exactly
\eqref{eq:elimination-schur}.  The inverse center records
have arbitrarily small derivatives in $\xi$ by that assertion;
the replay preserves them to any further fixed precision.  Hence
$|\sigma|\ge1/2$, all inverse powers are controlled, and the
initial calibration's record discrepancy yields a physical hit error
inside the fixed image neighborhood.  Correct the hit by
Lemma~\ref{lem:elimination-local-inverse}.

The center estimates, arbitrary replay accuracy, and bounded fixed
derivatives give $c_\ell=o(\rho^{m-\ell})$ at the actual Taylor
centers.  Therefore \eqref{eq:elimination-large-w} again holds.
This case immediately returns a visible output; it cannot occur for
exact zeros and does not pass an older numerical floor further down
the recursion.

\paragraph{The coarse case and its uniform conditioning powers.}
The only remaining case has center floor $\delta_A$, negligible $H$,
and negligible $w_A$.  Retain the center's active flag extension and
reference records.  Moving from $(y_A,w_A)$ to $(y_A,0)$ is smaller
than every floor power in bounded backgrounds, so it preserves the
floor comparisons.  Fix an eventual realization $r$, a field
tolerance $M$, and a derivative request $k$.

Choose a thickness exponent $J_0$ from $M$ and a center field
tolerance $M_H$ large enough for Taylor's formula to imply
\eqref{eq:elimination-coarse} when
\begin{equation}\label{eq:elimination-thickness}
 w=r^{J_0}u,
\end{equation}
with bounded $u$.  These choices do not involve record accuracy.
Explicitly, $H\le r^{M_H}$ implies
$|c_\ell|\le r^{M_H(m-\ell)/(2L)}$; choose each resulting
exponent larger than $M$ with a fixed slack, and choose $J_0m>M$
with slack.  Bounded Taylor coefficients and $a$ then give the
desired bound after a sufficiently accurate lift.

Apply the smaller-$n$ assertion, including its prefix invariant,
to preassign bounds for the center maps $Y^{[\mathcal A]}$, records $R^{[\mathcal A]}$, and
their derivatives through a sufficiently high fixed order, inverse
Jacobians, and source and image neighborhoods.  Their exponents do
not depend on the requested record accuracy $\mathcal A$.
The floor realization $r$ is a strict positive monomial times a bounded
regular unit with positive limiting value.  Lemma~\ref{lem:elimination-common-domain}
therefore applies to this center family with the already fixed new
thickness exponent $J_0$.
It gives a source exponent and regular-stage trial-tube exponents
before $\mathcal A$ is chosen.  On those same domains its replay can be
made arbitrarily accurate in the fixed differentiated norm,
increasing the finite Taylor depth only after the losses in the
selected flattened formula are known.  All new defining equations
and analytic arguments have been checked on genuine ordinary boxes
in that lemma, before positivity or a regular inverse is invoked.
Consequently the \emph{actual} lifted maps and records inherit the
preassigned power bounds on a preassigned power-sized domain.
Larger powers inside a finer formula neither reduce its available
source-radius exponent nor worsen these actual-map bounds.

Before lifting, the lower map is independent of $u$.  At fixed
physical $y$ the lifted inverse records thus have $u$-derivatives
$O(r^{T-C})$, by \eqref{eq:elimination-inverse-record}; their
exogenous derivatives differ from the arbitrarily small center
derivatives by an error of the same kind.  Write $L=\log r$ and
use the completed maps $Y(q,u),R(q,u),W(q,u)$, as in the lift lemma.
With $K_u=R_u-R_qY_q^{-1}Y_u$, their exact physical Schur complement is
\begin{equation}\label{eq:elimination-thickness-schur}
 W_u-W_qY_q^{-1}Y_u
 =r^{J_0}\frac{1+uJ_0L_RK_u}{1-wJ_0L_w}
 =r^{J_0}(1+o(1)).
\end{equation}
Here $L$ has no explicit dependence on the newly added $u$, and its
displayed derivatives hold physical variables and records fixed.
All implicit record feedback is already in $K_u$.  Its smallness and
the logarithmic sensitivity bounds give the last equality.  The full
inverse therefore loses only a fixed power determined by $J_0$ and
the center bounds.

The same fixed-power conclusion holds for the forward derivatives:
they are those of the replayed center maps and
\eqref{eq:elimination-thickness}, with the current-$w$ equation
initially solved by its regular unit Jacobian.  Record outputs are
either active free coordinate projections or positive ordinary
expressions governed by Theorem~\ref{thm:calculus}, so the required
bounds include their derivatives.  They are not arbitrary outputs
with uncontrolled growth.  Lemma~\ref{lem:elimination-local-inverse}
now supplies source and image radii with exponents fixed from these
bounds.  Thus all conditioning and neighborhood powers have been
chosen before the final record accuracy.

Given any requested $\mathcal A$, take the lower record accuracy and the replay
depth large enough to absorb these fixed inverse and output-derivative
losses.  At the desired physical point choose
$u=w_A/r^{J_0}$ using the reference records, as in the lift lemma.
This tends to zero because $w_A$ is negligible.  Calibration,
inverse correction, and the record derivative bounds give an exact
hit with record error $O(r^{\mathcal A})$.  Solving at fixed $(y,w)$ costs only
the same fixed powers.  More explicitly, for an exogenous coordinate
$e$ put $K_e=R_e-R_qY_q^{-1}Y_e$, using completed derivatives
at fixed $u$.  Differentiating the thickness identity at fixed
physical $(y,w)$ gives
\[
 \left.\partial_eu\right|_{y,w}
 =-\frac{uJ_0(L_e+L_RK_e)}{1+uJ_0L_RK_u},\qquad
 \left.\partial_eR\right|_{y,w}
 =K_e+K_u\left.\partial_eu\right|_{y,w}.
\]
The first expression has a fixed power bound, while $K_e,K_u$
can both be made smaller than any prescribed power on the common
domain.  Thus the second is $O(r^{\mathcal A})$ after requesting the lower
accuracies beyond that fixed loss.  Earlier completed charts have
no dependence on the
pending exogenous inputs: at fixed original physical input their
intermediate physical inputs are fixed as well.  Composing their
power-conditioned inverses preserves the required smallness.  This
proves the entire strengthened coarse assertion.

To make the prefix part explicit, the new structural list consists
of the fixed thickness equation and the smaller-dimensional list
replayed in its former order.  The limiting thickness equation is
$w=0$ with unit $w$-Jacobian, and each replayed regular equation or
analytic argument has its center's limiting germ.
Lemma~\ref{lem:elimination-common-domain} supplies its prefix
domains and finite derivative bounds.  The only new monomial
generators are the monomial factor of the fixed width $r^{J_0}$ and
the finitely many triangular factors needed when differentiating it; repeated
differentiation changes their powers.  Its inverse block is
controlled by \eqref{eq:elimination-thickness-schur}.
Thus the new list and all prefix exponents are fixed before the
record accuracy, as required.

\paragraph{Deferred floors and denominator accounting.}
There is only one outstanding numerical floor in a negligible
output.  A negligible center relative to visible $|w_A|$ realizes
that floor immediately and exits visibly.  A visible center chart
is lifted and its exact hit completed before further recursion.
Consequently neither kind of active decision depends on a future
realization of an inherited floor.  Along a chain which still passes
the inherited floor into centers, every removed $w_A$ is negligible
relative to it.  Old free scales remain fixed and bounded backgrounds
move by negligible powers.  Put $M=|\log\delta_A|$.  At a fixed
threshold $f_i=\log Z_i=O(M^q)$, the triangular relation for $f_i$
has largest participating child scale $O(M^q)$.  The finite triangular
differentiation used in the proof of
Lemma~\ref{lem:elimination-promotion} therefore bounds every fixed
background derivative by $O(M^D)$ for some fixed $D$.  This argument
uses the polynomial scale-log threshold directly; it does not require
the corresponding exponential monomial to be power moderate in
$\delta_A$.  Integration shows that a background displacement
negligible relative to every power of $\delta_A$ cannot carry a scale
log across a fixed comparison range about that threshold.  Thus scale
logs of at most a fixed power of $M$ retain their comparisons, and a
much larger log cannot enter their comparison range.  After realization
these are equivalently lexicographic monomial comparisons in the
extended saturated flag.

Reference records of completed charts are taken at exact hits.
Those of a pending center output are taken at its center projection;
negligible center displacements and the eventual arbitrary
power-accurate hit preserve their comparisons.  Pending floor
recipes add only thickness parameters replacing consumed ordinary
coordinates, not active promotions tied to an extra recursion
variable.  Thus the induction really decreases $n$ or, at fixed $n$,
decreases $m$.

Finally keep a list of inverted monomials at each step.  Outside an
immediate negligible exit the leading $a$ is moderate.  The
multiplicity step adds only $s$ and its fixed powers, moderate by
visibility and the inequalities already proved.  Coarse outputs
use only floor-moderate scales.  The flattening and all fixed
derivative formulas introduce finitely many further powers of these
same denominators.  Each recipe is therefore power moderate in its
required scale.  The family assertion is stronger only for actual
maps and their neighborhoods, whose uniform exponents were proved
by finite differentiability comparison above.  This completes both
inductions.
\end{proof}

\subsection{The absolute zero theorem and its closure}

The final retest requires a property of the primitive library, in
addition to the estimates used in Theorem~\ref{thm:calculus}.  We
state it explicitly to separate a local formal identity from a
claim about arbitrary changes of asymptotic regimes.

\begin{definition}\label{def:elimination-retestable}
A packet library is \emph{retestable for elimination} if it satisfies
Theorem~\ref{thm:calculus} on each saturated extension used above and
has the following germ invariance.  Arbitrarily small absolute
perturbations of retained old free inputs which preserve the affine
flag relations, regime gaps, limiting ordinary data, and fixed sector
determinations can be tested with the same terminal ordinary formal
coefficient germs.  Refining by unused nodes inserts identity
transitions before pullback.  Numerical anchor and cutoff choices
must disappear by the specified formal coefficient recursions or
normalized limiting jets before terminal coefficients are reached.
The property concerns the selected local lifted germs and does not
assert single-valuedness over all free-input histories.
\end{definition}

\begin{theorem}[Absolute isolated-zero finiteness]
\label{thm:absolute-zero}
Let a finite real equation system be specified on an open domain
with its original strict inequalities and graph ties.  Regard all
coordinates, including coefficient parameters and auxiliary
coordinates, as variables.  Suppose every sequence of distinct
solutions has a subsequence with one of the following descriptions:
\begin{enumerate}
\item the equations extend as ordinary real analytic functions to a
neighborhood of its finite limit; or
\item after finitely many regular graph constructions and coordinate
charts, it is an anchored saturated flag regime for a library
retestable in the sense of Definition~\ref{def:elimination-retestable},
with at least one old free input recoverable from the original
solution tuple.
\end{enumerate}
In the second description every equation admits the finite
normalizations and positive operations of
Theorem~\ref{thm:calculus}.  Then the system has finitely many isolated
real solutions in its full domain.

The assertion applies separately to every finite system in the
following closure whenever these same sequential hypotheses hold:
independent derivatives of equations and domain margins; finite sums,
products and regular analytic evaluations; inverses on their specified
nonzero domains; denominator clearing with those domains retained;
fixed real specializations; addition of real parameters, multiplier
variables, polynomial equations, and regular graph ties, including
finite derivative graphs.  No uniform bound over infinitely many
generated systems, and no assertion after forgetting their original
domains, is part of this theorem.
\end{theorem}

\begin{proof}
Suppose there is an infinite sequence of distinct isolated solutions.
In the first alternative a real analytic zero set has locally finitely
many connected components \cite[Corollary~2.7]{BierstoneMilman1988},
so its isolated points cannot accumulate at an interior point of an
analytic extension neighborhood.  Intersecting with an open original
domain cannot turn
a positive-dimensional local analytic branch at one of its interior
points into an isolated point.

In the second alternative, multiply each equation by an appropriate
nonvanishing monomial so that it is positive, and replace the finite
system by the sum of the squares of these real normalized equations.
The zero set on the retained domain is unchanged.  The sum is
positive by the calculus.  Apply exact mode of
Proposition~\ref{prop:elimination-recursion}.  At every sufficiently
far anchor the resulting chart satisfies identically zero formal
data.  Ordinary restrictions may have reduced its ordinary dimension,
but they have never removed an old independent large input.

Take an isolation neighborhood about each alleged isolated solution.
In its final chart choose a nonzero absolute perturbation of a
recoverable old free input small enough that the mapped point remains
inside that neighborhood, inside all strict domains and chart
neighborhoods, and in the same saturated regime.  The perturbations
can be successively smaller than every finite list of required
bounds; they preserve the limiting ordinary data and exact affine
flag ties.  Recoverability makes their solution tuples different
from the original ones.

Retest on this perturbed sequence.  The primitive terminal germs are
the same by Definition~\ref{def:elimination-retestable}.  Every later
formal operation also gives the same germ: algebra and derivatives
act on the same ordinary identities, finite Taylor replays use the
same formulas, and regular implicit equations select the same
ordinary root germs.  In particular, both all-zero formal arrays
remain zero.  Separation therefore says that the actual normalized
equations vanish at all sufficiently far perturbed points.  These
are distinct solutions in the proposed isolation neighborhoods, a
contradiction.

The closure statement means applying this argument to each particular
finite generated system.  Every allowed operation is a finite
calculus operation or a regular graph construction; polynomial
relations in newly added variables are ordinary analytic relations.
A derivative or multiplier system must still meet the sequential
hypotheses, with every original strict domain and graph tie retained.
Clearing a denominator on its nonzero domain preserves its zero set.
Thus the argument applies to that system with all its coordinates
variable.  It does not replace this hypothesis by finiteness of the
original unaugmented equation set.
\end{proof}

\begin{remark}
The concrete primitive constructions below verify the retest property
by coefficient recursions, compatible finite-charge quotients, or
canonically normalized infinite-anchor jets.  Merely discarding an
unknown flat difference would not verify the equality of terminal
germs needed in the last proof.  The sequence-dependent flags and
finite internal recipes used here likewise do not assert a uniform
number of those recipes before a sequence is selected.
\end{remark}
